\section{引言}\label{sec:intro}

\subsection{问题的提出}

图上的很多任务，人类已经写出了“最优解”：单源最短路有 Dijkstra 与 Bellman-Ford，最小生成树有 Kruskal 与 Prim，最大流有 Dinic。
这类算法的共同点是\textbf{白盒}的——每一步都可以被审计，正确性有证明，复杂度有上界，输入被扰动之后答案会怎么变，往往也能算出来。

与之相对，图神经网络（GNN）走的是另一条路：把图上的计算交给一层层消息传递，用数据拟合出映射。
它在节点分类、链接预测上常常表现很好，但\emph{为什么}给出这个答案，只能靠事后解释器去近似~\cite{ying2019gnnexplainer}。
这类模型是\textbf{黑盒}的：假设类无法逐条读出，误差没有上界，解释是近似的近似。

于是题目给出的那个“vs.”就成立了：\textbf{可解释性（Interpretability）vs. 性能（Performance）}。
但把它写成一句口号很容易，量化却很难：什么叫可解释？鲁棒性怎么度量？准确率用什么口径？
本文的做法是\emph{不并列介绍两类方法}，而是把它们放进同一个任务、同一批数据、同一套指标里去比。

\subsection{为什么选最短路作为载体}

选择单源最短路有三个理由。

第一，它有\textbf{精确解}。Dijkstra 与 Bellman-Ford 的正确性可以证明，因此白盒一侧的“准确率”不是估计出来的，而是恒等于真值；任何黑盒的误差都可以被精确地量出来，而不是和另一个估计比。

第二，它的\textbf{结构与消息传递同构}。Bellman-Ford 的第 $k$ 轮只做一件事：沿边松弛一次。而 $k$ 层消息传递的感受野恰好是 $k$ 跳。两者的对应不是比喻，是可以逐项写下来的（定理~\ref{thm:equiv}）。这条桥一旦搭起来，“深度”“聚合方式”“可解释性”就都能落到具体的位置上。

第三，它\textbf{可以完全在 CPU 上复现}。本文所有实验只用合成图（ER / BA / 网格 / 随机块模型），不依赖任何外部数据集，单个脚本在几十秒内跑完，任何人都能重跑。

\subsection{本文的工作}

\begin{enumerate}[leftmargin=1.8em]
\item \textbf{理论}：证明 $k$ 层 min-sum 消息传递与同步 Bellman-Ford 的第 $k$ 轮逐点等价（定理~\ref{thm:equiv}）；
      给出深度下界的构造性反例（命题~\ref{prop:depth}），说明 $L$ 层模型在两张 $L$ 跳计算图同构、但真值不同的图上必然出错，且这个错误\emph{与训练无关}。
\item \textbf{白盒侧}：把“鲁棒性”从形容词变成三个可证命题——删边单调性、非树边不变性、单边删除的可计算上界（命题~\ref{prop:mono}--\ref{prop:bound}），并在 500 个随机实例上逐例验证，零违反。
\item \textbf{统一实证}：设三档模型（L0 恒等 min-sum、L1 训练后的 2 参数 min-sum、L2 约 $10^5$ 参数的 mean-GNN），在准确率、鲁棒性、可解释性三轴上测量。
      三个反直觉的结果是：
      \begin{itemize}[leftmargin=1.4em]
      \item \textbf{训练是负收益}：L0 的误差是 $\sim10^{-6}$（float32 舍入级），训练后的 L1 是 $0.024$——训练把参数推离了精确解，而迭代会把这点偏差放大；
      \item \textbf{“不敏感”不等于“稳健”}：黑盒的误差几乎不随权重扰动变化，因为它的误差本来就已经接近常数预测器的水平；
      \item \textbf{保真度会被常数预测器骗过}：黑盒的 Fidelity$^+$ 是 0.978，看起来解释得很好，但它的预测与真值的相关系数只有 0.13。
      \end{itemize}
\end{enumerate}

\subsection{全文结构}

第~\ref{sec:white}~节给出白盒侧的算法与轮数语义；第~\ref{sec:robust}~节把鲁棒性写成可证命题并验证；
第~\ref{sec:black}~节建立消息传递与 Bellman-Ford 的等价，并给出深度下界；第~\ref{sec:setup}~节说明实验设计；
第~\ref{sec:results}~节报告三轴结果；第~\ref{sec:discussion}~节讨论结论的边界；第~\ref{sec:conclusion}~节总结。
